% =====================================================================
%  Chapter 13 — The YORP \idxw{Torque}
% =====================================================================

\chapter{The YORP Torque}
\label{ch:yorp}

The YORP effect --- Yarkovsky--O'Keefe--Radzievskii, named for the three
authors whose separate predictions of it turned out to describe the same
phenomenon --- is the torque\idx{Torque} that a rotating body with a non-uniform surface
exerts on itself through the anisotropic re-radiation of sunlight. It is the
mechanism by which a small asteroid's rotation rate changes over decades,
and it is the channel through which a shape measurement does its most
direct work on the probability chain.

\section{The mechanism}

Consider a surface element of a rotating body. It is heated as it rotates
into sunlight and cools as it rotates away. The temperature of that element
therefore varies through the rotation, and the radiation it emits varies with
it. The emitted radiation carries momentum, and the momentum is directed
along the outward normal. Integrating over a surface whose normals are not
uniformly distributed leaves a net torque.

Three properties of the surface determine the result: the distribution of
surface normals, the distribution of roughness, and the obliquity. The
first is what a shape model supplies, and it is why the effect is
shape-sensitive in a way that the Yarkovsky force is not.

The published scaling, from the same 2012 study, is
\begin{equation}
\frac{d\omega}{d t} \propto \frac{1}{\rho_b D^{2}},
\label{eq:yorp}
\end{equation}
with the dependence on the shape and the roughness carried separately. The
diameter enters as the square, and the density linearly.

\begin{table}[!htbp]
\centering
\small
\caption{What the YORP timescale means for the small class. The calibration
object has a torque of \(2.26\times10^{-3}\) radians per year squared.}
\label{tab:yorp-timescale}
\begin{tabularx}{\textwidth}{@{}P{0.91in}P{0.88in}P{0.91in}Y@{}}
\toprule
\rowcolor{panel}
\textbf{Diameter} & \textbf{Density} & \textbf{Torque ratio} &
\textbf{Comment} \\
\midrule
51 m & 1,300 & \(\times 4{,}827\) & the torque\idx{Torque} on this class is
enormous relative to the calibration \\
\rowcolor{panel}
60.8 m & 1,300 & \(\times 3{,}409\) & the YR4 class \\
60.8 m & 2,800 & \(\times 1{,}583\) & as above, coherent \\
\rowcolor{panel}
140 m & 2,800 & \(\times 299\) & the catalogue's threshold class \\
\bottomrule
\end{tabularx}
\end{table}

The ratios in \tabref{tab:yorp-timescale} are the size-and-density part
only, and they are large. A torque\idx{Torque} that would take a 2.5-kilometre object a
decade to spin up would spin a 60-metre object of the same density in a
fraction of a year --- which is physically absurd, and is the reason the
same caveat applies here as in the previous chapter: the full YORP
expression carries a dependence on the shape and on the spin rate that does
not scale with the diameter alone, and the book does not extrapolate an
absolute secular acceleration across two orders of magnitude.

What survives the caveat is the comparative statement, and it is the one
that matters: \textbf{the YORP torque on the small class is enormously larger
relative to its calibration than the Yarkovsky drift is}, by a factor of
roughly fifty in the ratios of \tabref{tab:yorp-timescale} against
\tabref{tab:rg-ratios}. For the small class, the rotation state is a fast
variable, and a fast variable is a variable that a characterisation
campaign has to catch rather than assume.

\section{Why this matters for the book}

Three consequences follow, and they are the reason this chapter exists.

The first is that the spin state is not a slowly varying property that can
be carried forward from a survey lightcurve. For this class it evolves on a
timescale shorter than the gap between observing opportunities. A catalogue
that stores a rotation period measured once is storing a quantity with a
shelf life.

The second is that the shape enters the probability chain directly. The
published sensitivity result is that a three per cent error in the diameter
produces a six per cent error in the YORP torque --- twice the relative
sensitivity of Yarkovsky --- because the torque depends on the shape
directly and the drift does not.

The third is the one that connects this chapter to Chapter 7 and to the
whole architecture. A torque is a vector, and the direction it turns the
body towards depends on where the mass is. The 2022 study of Apophis's
rotation through the 2029 encounter found that the encounter is likely to
change the effective spin rate by up to a factor of two and to shift the
pole by ten degrees or more, and that these changes are very sensitive to
the close-approach attitude and the mass distribution. The spin state is
therefore not merely an input to the torque; it is the \emph{output} of the
encounter, and it is the observable through which the mass distribution
becomes measurable.

\begin{seamnote}{The chain closes in a place nobody expected}
The sequence is: the encounter perturbs the rotation; the change in
rotation is measured from the ground by photometry and a resolved shape; the
change in rotation constrains the mass distribution; the mass distribution
gives the density; the density is the second factor of three in the energy
relation. Every one of these steps is a measurement that a station at the
encounter could support, and the last one is the one the catalogue has never
had.
\end{seamnote}

\section{What a lightcurve alone can and cannot do}

The lightcurve is the workhorse here, and it is worth being precise about
what it delivers and where it stops.

It delivers the rotation period, the direction of the spin pole, and a
low-order shape model, and it does so from a telescope aperture that every
major observatory already owns. It is cheap, it is fast, and it does not
require proximity.

It stops at the resolution of the shape. A lightcurve constrains the
projected outline of the body, and the standard inversion returns the convex
hull of that constraint. The surface normals that the YORP torque depends on
are the normals of the body, not of its convex hull, and the difference is
exactly the concavity that the hull cannot represent.

For a body in this class, with a spin rate that may change by a factor of
two across a close encounter, the shape resolution is the binding limit on
the torque. This is not a limitation of the technique; it is a statement
about what the technique measures. Chapter 17 takes the convex-hull problem
apart in detail, and Chapter 20 works out what an interferometric
measurement of the same body would deliver.

\begin{keytakeaways}
\item YORP depends on the surface normals directly, which is why it is
genuinely shape-sensitive where Yarkovsky is not.
\item The published scaling puts the torque on the small class tens to
hundreds of times the calibration relative to its size and density, and the
book does not extrapolate an absolute secular acceleration across that
range.
\item A three per cent diameter error produces a six per cent YORP error,
twice the relative sensitivity of Yarkovsky.
\item The spin state is the output of a close encounter as well as its input,
and it is the observable through which the mass distribution becomes
measurable.
\end{keytakeaways}
